\chapter*{Wstęp}
\addcontentsline{toc}{chapter}{Wstęp}

Rozwój mierników mocy, czujników tętna i platform treningowych umożliwia systematyczne gromadzenie danych z codziennych aktywności kolarskich. Dane te mogą uzupełniać okresowe testy wydolnościowe, które są obciążające i trudne do częstego powtarzania. Problem badawczy podjęty w pracy dotyczy możliwości estymacji Functional Threshold Power (FTP) na podstawie zagregowanej telemetrii z rutynowych jazd, bez wykorzystywania cech bezpośrednio odtwarzających sposób wyznaczenia wartości referencyjnej.

FTP jest w pracy traktowane jako operacyjny wskaźnik dyspozycji wytrzymałościowej, a nie pełna miara formy sportowej ani odpowiednik progu zmierzonego laboratoryjnie. Wartość referencyjną stanowi heurystyczna etykieta obliczona jako 95\% maksymalnej średniej mocy z 20 minut. Przyjęcie takiej definicji umożliwia ilościową ocenę modeli, ale ogranicza zakres interpretacji wyników.

Sformułowano następującą hipotezę badawczą:
\begin{quote}
    \textit{Co najmniej jeden model regresyjny wykorzystujący zagregowane cechy telemetryczne, z wyłączeniem całej rodziny cech Maximum Mean Power (\texttt{mmp\_*}), osiągnie na zbiorze testowym obejmującym zawodników nieobecnych w zbiorze uczącym błąd MAE nie większy niż 15 W oraz współczynnik determinacji $R^2$ nie mniejszy niż 0{,}80 w estymacji operacyjnej etykiety FTP.}
\end{quote}

Celem głównym pracy jest zaprojektowanie, implementacja i ewaluacja prototypowego systemu analitycznego do estymacji operacyjnej etykiety FTP. Cel poznawczy obejmuje usystematyzowanie wiedzy o FTP i analizie telemetrii kolarskiej. Celem metodologicznym jest przygotowanie powtarzalnego potoku przetwarzania danych i rzetelnej procedury oceny modeli z kontrolą data leakage. Cel utylitarny polega na ocenie, czy wynik modelu może stanowić dodatkową informację w monitorowaniu zmian FTP.

Badanie odpowiada na następujące pytania:
\begin{enumerate}
    \item Czy modele uczone bez cech \texttt{mmp\_*} spełniają przyjęte kryteria jakości na danych zawodników nieobecnych w zbiorze uczącym?
    \item Który z algorytmów: Ridge Regression, Random Forest i XGBoost, osiąga najlepsze wartości MAE, RMSE i $R^2$?
    \item Jak zmienia się jakość predykcji po usunięciu cechy \texttt{mmp\_20min}, a następnie całej rodziny cech \texttt{mmp\_*}?
    \item Które grupy cech mają największe znaczenie dla predykcji?
\end{enumerate}

Zakres przedmiotowy obejmuje przetwarzanie telemetrii kolarskiej, inżynierię cech i regresyjną estymację FTP. Zakres podmiotowy tworzą anonimowi użytkownicy repozytorium GoldenCheetah OpenData. Zbiór bazowy przygotowany z 781 archiwów ZIP obejmuje po filtracji 151\,399 sesji 726 zawodników. Zakres czasowy wyznaczają historyczne aktywności dostępne w wersji repozytorium pozyskanej do badania, a zakres przestrzenny ma charakter międzynarodowy i wynika ze struktury tego repozytorium.

Zastosowano analizę literatury, przygotowanie potoku Extract--Transform--Load, inżynierię cech oraz eksperyment porównawczy modeli Ridge Regression, Random Forest i XGBoost. Dane podzielono na poziomie zawodników z użyciem \texttt{GroupShuffleSplit}, a jakość oceniono metrykami MAE, RMSE, MedAE i $R^2$. Wariant kontrolny A zawierał \texttt{mmp\_20min}, wariant B wykluczał tę cechę, a wariant C wykluczał wszystkie cechy \texttt{mmp\_*}.

Praca składa się z czterech rozdziałów merytorycznych. Rozdział pierwszy przedstawia podstawy teoretyczne i przegląd literatury. Rozdział drugi opisuje metodologię oraz projekt rozwiązania, trzeci prezentuje eksperymenty i wyniki, a rozdział czwarty zawiera ocenę rozwiązania, ograniczenia oraz możliwości praktycznego zastosowania wyników. Zakończenie podsumowuje wyniki i weryfikuje hipotezę.

Podstawę literaturową stanowią publikacje dotyczące treningu z pomiarem mocy, wiarygodności FTP, fizjologii wysiłku oraz zastosowań uczenia maszynowego i zasad rzetelnej walidacji modeli \parencite{allen2019power,borszcz2018ftp,karsten2021cpftp,stockwell2023mlftp,denham2020predictftp,james2021isl,hastie2009esl}. Dane empiryczne pochodzą z repozytorium GoldenCheetah OpenData \parencite{osf_goldencheetah_opendata}.
